\subsection{Propositions 7--9 and the gluing axioms}

\begin{definition}[Known: the full gluing axioms of XII]\label{A:audit-gluing}
For a fundamental localizer $W$, the axiom called W(7 bis) at XII ms.~75--76 asserts the following biconditional for every small category $X$ covered by two full sieves $U,V$ \citep{GrothendieckDerivateursXIIb}:
\[
 (U\cap V\longrightarrow V)\in W
 \quad\Longleftrightarrow\quad
 (U\longrightarrow X)\in W.
\]
The symmetric biconditional is included. W(7), at ms.~76--77, concerns a morphism of two such cover squares. Of the three conditions that the map on the intersection belongs to $W$, that both maps on the pieces belong to $W$, and that the map on the union belongs to $W$, any two must imply the third. These are stronger assertions than forward homotopy gluing.
\end{definition}

\begin{lemma}[Proved: a categorical model for a one-point attachment]\label{A:audit-cylinder}
Let $U$ be a small category, $H:U^{\mathrm{op}}\to\mathbf{Set}$ a presheaf, and $E=U/H$ with projection $\alpha:E\to U$. Let $X$ be the category with full sieve $U$, one complementary object $t$ with only its identity endomorphism, and crossing arrows $\operatorname{Hom}_X(u,t)=H(u)$. There is a category $M$, covered by two full sieves $M_U,M_t$ with intersection $E$, and a functor $F:M\to X$, such that
\[
 r:M_U\longrightarrow U\ \in W_\infty,
 \qquad M_t=E^{\triangleright},\qquad F\in W_\infty.
\]
The map $Br$ is a homotopy equivalence induced by a specified adjunction. Consequently $BX$ is the homotopy pushout of $BU\leftarrow BE\to *$.
\end{lemma}

\begin{proof}
Take disjoint copies of $E$ and $U$ and a new object $t$. Retain all their internal arrows, set
\[
 \operatorname{Hom}_M(e,u)=\operatorname{Hom}_U(\alpha e,u),
 \qquad \operatorname{Hom}_M(e,t)=\{*\},
\]
and allow no other crossing arrows. Composition uses composition in $U$ and the unique arrows into $t$. Put $M_U=E\cup U$ and $M_t=E\cup\{t\}$. Both are sieves and their intersection is exactly $E$. The functor $r:M_U\to U$ sending $e$ to $\alpha e$ and fixing $U$ is left adjoint to the full inclusion $U\to M_U$: the required Hom-set bijections are the defining formula above. Its counit is the identity and its unit is given by $e\to\alpha e$, so $Br$ is a homotopy equivalence.

Define $F:M\to X$ by $F(e)=\alpha e$, $F(u)=u$, and $F(t)=t$. For $e=(u,h)$, send its unique arrow to $t$ to $h\in H(u)$. Every other crossing arrow is sent to its defining arrow of $U$. The relation $H(g)(k)=h$ defining an arrow $(u,h)\to(v,k)$ in $E$ proves functoriality at $t$.

For $u\in U$, the comma category $M/u$ has terminal object $(u,\mathrm{id}_u)$. The only remaining comma category is $M/t$. Its $E$-layer objects are triples $(u,h,k)$ with $h,k\in H(u)$: $(u,h)$ is the object of $E$ and $k:F(u,h)\to t$ is the comma arrow. Its $U$-layer objects are $(u,k)$ with $k\in H(u)$, and there is one further object $(t,\mathrm{id}_t)$. An $E$-layer object has an arrow to the last object precisely when $h=k$.

Let $N\subseteq M/t$ contain the $U$-layer, the last object, and just the diagonal triples $(u,h,h)$ in the $E$-layer. Define $S:M/t\to N$ to fix $N$ and send an off-diagonal $(u,h,k)$ to the $U$-layer object $(u,k)$. Send each arrow to the arrow with the same underlying morphism of $U$, or to the unique arrow into the last object. This prescription is defined on every arrow: an arrow of $E$-layers whose target is diagonal has diagonal source, since restriction of equal sections gives equal sections. There is a natural transformation $\mathrm{Id}_{M/t}\Rightarrow\iota S$ whose only nonidentity components are $(u,h,k)\to(u,k)$ with underlying arrow $\mathrm{id}_u$. Thus $N\hookrightarrow M/t$ is a homotopy equivalence.

The category $N$ consists of two copies of $E$, the first mapping to the second by the categorical mapping cylinder of the identity, and arrows from the first copy to $t$. The endofunctor $R:N\to N$ sending the second copy to the first and fixing the first copy and $t$ has natural transformations
\[
 \mathrm{Id}_N\ \Longleftarrow\ R\ \Longrightarrow\ \mathrm{const}_t.
\]
The first uses the cylinder arrows, and the second uses the unique arrows into $t$. Hence $N$, and therefore $M/t$, is contractible. Quillen's Theorem A, in its overcategory form \citep[\S1, Theorem~A]{Quillen1973}, proves $F\in W_\infty$.

Every nerve simplex of $M$ lies in $M_U$ or $M_t$, because these are a sieve cover. Their nerve inclusions are monomorphisms of simplicial sets and their realizations are cofibrations. The literal union $BM=BM_U\cup_{BE}BM_t$ is therefore a homotopy pushout. Applying $BF$ and $Br$ proves the last assertion, with its specified comparison maps. The argument includes $H=\varnothing$ and $U=\varnothing$.
\end{proof}

\begin{proposition}[Proved: Proposition 7 with its stated additional axiom]\label{A:audit-prop7}
Let $W$ be a fundamental localizer satisfying the full W(7 bis), and let $U\subseteq X$ be a full sieve. Write $T=X\setminus U$, and let $\overline U$ be the full cosieve generated by $U$. The following are equivalent: the inclusion $U\to\overline U$ is $W$-aspherical; for every functor $g:T_1\to T$, replacing the complement by $T_1$ and the crossing bifunctor by $(u,s)\mapsto\operatorname{Hom}_X(u,g(s))$ gives an inclusion $U\to\overline U_1$ in $W$; and this last condition holds for every point $g:e\to T$ whose value belongs to $\overline U\setminus U$.

For a one-point complement, the corollary at XII ms.~85 is likewise valid under full W(7 bis): $U\to X$ belongs to $W$ if and only if $U/t$ is $W$-aspherical.
\end{proposition}

\begin{proof}
We first prove the one-point assertion. Use Lemma~\ref{A:audit-cylinder} for $H(u)=\operatorname{Hom}_X(u,t)$, so $E=U/t$. The two functors $r:M_U\to U$ and $F:M\to X$ belong to every fundamental localizer, since they belong to $W_\infty$. By two-out-of-three and the equality $F|_{M_U}=ir$, the map $M_U\to M$ belongs to $W$ exactly when $i:U\to X$ does. W(7 bis), applied to the sieve cover $M=M_U\cup M_t$, identifies this with $E\to M_t$ belonging to $W$. Since $M_t$ has terminal object $t$, two-out-of-three identifies the latter condition with $E\to e$ belonging to $W$.

For the three conditions in the proposition, an object $s$ of the replaced complement has comma category $U/s=U/g(s)$, literally with the same objects, arrows and composition. Objects of $U$ have comma categories with terminal objects. Thus the first condition implies $W$-asphericity, hence membership in $W$, of each replaced inclusion by the relative Theorem A axiom. The second condition includes the third. The one-point assertion shows that the third condition makes $U/t$ $W$-aspherical for every $t\in\overline U\setminus U$, which is the first. If this set is empty, all point tests and the comma condition are vacuous and $U=\overline U$; thus no nonemptiness hypothesis has been suppressed.
\end{proof}

\begin{example}[Proved: an explicit finite acyclic category with nontrivial fundamental group]\label{A:audit-acyclic}
There is a category $A$ with $35$ objects and $159$ arrows, including identities, such that $BA$ is connected and integrally acyclic but has nontrivial fundamental group. We specify its complete finite presentation through face operators.

Start with the two words
\[
 w_0=aaaaaBBB,\qquad w_1=bbbBABA,
\]
where capital letters denote inverses, with lengths $n_0=8,n_1=7$. A finite semisimplicial set $K$ has vertices $v,c_0,c_1$, edges $a,b$ with both endpoints $v$, and edges $r_{k,i}:c_k\to v$ for $k\in\{0,1\}$ and $i\in\mathbb Z/n_k$. It has triangles $t_{k,i}$ with ordered vertices $(c_k,v,v)$. If the $i$th letter of $w_k$ is the positive letter $\ell$, define
\[
 (d_0t_{k,i},d_1t_{k,i},d_2t_{k,i})
 =(\ell,r_{k,i+1},r_{k,i});
\]
if that letter is $\ell^{-1}$, define this triple to be $(\ell,r_{k,i},r_{k,i+1})$. The endpoint faces of every $r_{k,i}$ are $(v,c_k)$, and those of $a,b$ are $(v,v)$. These prescriptions satisfy the semisimplicial identities. There are no simplices above dimension two.

The objects of $A$ are the simplices just listed. A morphism $\sigma\to\tau$ is an increasing injection $\theta:[\dim\sigma]\hookrightarrow[\dim\tau]$ for which $\theta^*\tau=\sigma$; composition is composition of the injections. In particular, distinct occurrences of the same face give distinct arrows. These definitions specify every object, arrow and composite of $A$.
\end{example}

\begin{proof}
The realization $|NA|$ is canonically homeomorphic to the barycentric subdivision of $|K|$. Indeed, each nondegenerate nerve simplex is a flag of proper face injections; the barycentres of the nested faces specify its usual barycentric simplex in the final target simplex. These simplexes subdivide each characteristic simplex, and the constructions agree on every face identification, including repeated faces.

Geometrically $K$ triangulates the presentation complex with one vertex, two loop edges and two two-cells attached by $a^5b^{-3}$ and $b^3(ab)^{-2}$. For each attaching polygon the triangles above are its cone triangulation, with centre $c_k$ and radial edges $r_{k,i}$; reversing the order of its last two triangle vertices for a negative letter gives exactly the indicated attaching map. This is the acyclic presentation complex of \citet[Example~2.38, p.~142]{Hatcher2002}, now with explicit finite categorical data.

The cellular differential from degree two to degree one is
\[
 d_2=\begin{pmatrix}5&-2\\-3&1\end{pmatrix},
 \qquad \det(d_2)=-1.
\]
The complex is connected, $d_1=0$, and there are no higher cells. Hence its reduced integral homology vanishes. To verify nontrivial fundamental group, send
\[
 a\longmapsto(1\ 4\ 3\ 2\ 5),\qquad
 b\longmapsto(1\ 2\ 3)
\]
in $A_5$, using right-to-left permutation composition. Both relator words have identity image: the permutations have orders five and three, and $ab=(1\ 5)(3\ 4)$. Since the image of $a$ is not the identity, the presented group, and therefore $\pi_1(BA)$, is nontrivial.

The program \texttt{code/ag\_finite\_examples.py} generates all arrows and composites from these face operators and checks all category axioms. The normalized nerve chain groups have ranks $(35,124,90)$; the two boundary matrices have ranks $34,90$ and all their nonzero Smith factors are $1$. Thus the independent integral calculation gives $H_0=\mathbb Z$ and $H_i=0$ for $i>0$. It also checks the two relator permutations. Full arrow and composition tables are in \texttt{code/ag\_category\_data.json}.
\end{proof}

\begin{theorem}[Refuted: the unqualified one-point converse and W(7 bis) for $W_\infty$]\label{A:audit-prop7counter}
The converse ``$U\to X$ is a weak equivalence, therefore $U/t$ is weakly contractible'' fails for a finite sieve inclusion with a one-object complement. In particular $W_\infty$ satisfies neither the full W(7 bis) nor the full W(7), and Proposition~7 cannot be read with the additional axiom deleted.
\end{theorem}

\begin{proof}
Take $A$ from Example~\ref{A:audit-acyclic}. Adjoin two distinct terminal cone points $u,t$ over $A$, with no arrows between $u$ and $t$, and call the resulting category $X$. Thus every object $a\in A$ has one arrow to each cone point; composing any arrow of $A$ with such an arrow gives the unique corresponding arrow. Let $U=A\cup\{u\}$ and $V=A\cup\{t\}$. Both are full sieves, $X=U\cup V$, $U\cap V=A$, and $U/t\cong A$.

The nerves of $U$ and $V$ are cones on $NA$, so $BU$ and $BV$ are contractible and $BX\cong\Sigma BA$. Since $BA$ is connected, van Kampen gives $\pi_1(BX)=1$. The reduced Mayer--Vietoris sequence identifies $\widetilde H_i(BX;\mathbb Z)$ with $\widetilde H_{i-1}(BA;\mathbb Z)$, which vanishes. The simply connected acyclic CW complex $BX$ is contractible by the homological Whitehead theorem \citep[Corollary~4.33, p.~367]{Hatcher2002}. Hence $U\to X$ belongs to $W_\infty$, whereas $U/t\cong A$ is not $W_\infty$-aspherical. The cover square has $U\to X$ in $W_\infty$ and $A\to V$ outside $W_\infty$, contradicting W(7 bis). Mapping this cover square to the constant point square gives a direct failure of W(7): the two piece maps and the union map are weak equivalences, while the intersection map is not.

The category $X$ has $37$ objects and $231$ arrows. The same program checks its composition, both sieve conditions, and its integral homology. The topological argument, including the nontrivial fundamental-group certificate for $BA$, supplies the homotopy information that homology computation alone cannot establish.
\end{proof}

\begin{proposition}[Proved: the homology localizers satisfy the full gluing axioms]\label{A:audit-homologygluing}
For every commutative coefficient ring $R$, the class $W_R$ of nerve $R$-homology equivalences satisfies W(7), and consequently W(7 bis).
\end{proposition}

\begin{proof}
For a sieve cover $X=U\cup V$, any chain of arrows in $X$ lies in a piece containing its last vertex. Thus the normalized nerve chains give a natural short exact sequence
\[
 0\longrightarrow C_*(N(U\cap V);R)
 \xrightarrow{(i_*,-j_*)}C_*(NU;R)\oplus C_*(NV;R)
 \longrightarrow C_*(NX;R)\longrightarrow0.
\]
Exactness is checked on the bases of nondegenerate chains. For a morphism of cover squares, the mapping cones of the vertical chain maps again form a short exact sequence of complexes. The middle cone is the direct sum of the cones for the two piece maps. The long exact homology sequence proves that acyclicity of any two of the intersection cone, the pair of piece cones, and the union cone implies acyclicity of the third. A chain map is a homology equivalence exactly when its cone is acyclic, proving W(7). To obtain W(7 bis), map the given cover square to the cover of $X$ by $X$ and $V$. The maps on $V$ and $X$ are identities; W(7) then identifies the equivalence condition on $U\cap V\to V$ with that on $U\to X$.
\end{proof}

\begin{proposition}[Proved: Propositions 8 and 9(a)]\label{A:audit-prop8and9a}
For every fundamental localizer $W$, if $U/H$ is totally $W$-aspherical, then $H$ is totally $W$-aspherical. Every small category with a terminal object is totally $W$-aspherical, and every representable presheaf is totally $W$-aspherical.
\end{proposition}

\begin{proof}
For $\alpha:U/H\to U$, there is a canonical isomorphism $\alpha_!1_{U/H}\cong H$. Consequently, for every $p:U\to U'$, one has $p_!H\cong(p\alpha)_!1_{U/H}$; the first assertion is exactly the defining test for $U/H$ applied to $p\alpha$.

If $v_0$ is terminal in $V$ and $f:V\to B$, then
\[
 f_!1_V\cong B(-,f(v_0)).
\]
For an explicit verification, any object $(v,\beta:b\to f(v))$ of $b\backslash V$ maps to $(v_0,f(v\to v_0)\beta)$. This gives a retraction onto the discrete set $B(b,f(v_0))$, with a natural transformation from the identity to the inclusion followed by the retraction. Thus the stated component bijection is natural in $b$. Its category of elements is $B/f(v_0)$, with terminal object $(f(v_0),\mathrm{id})$, so the terminal-object axiom for $W$ applies. Finally $U/U(-,u)\cong U/u$ has a terminal object, and the first assertion gives the representable case.
\end{proof}

\begin{theorem}[Refuted: Proposition 9(b), its retraction repair, and the fork corollary]\label{A:audit-prop9b}
For $W=W_\infty$, and also for $W=W_R$ with any nonzero commutative ring $R$, a union of two totally $W$-aspherical full sieves with totally $W$-aspherical intersection need not be totally $W$-aspherical, even if both pieces retract onto the intersection. In the homology case this refutes the statements under their full stated W(7 bis) hypothesis.
\end{theorem}

\begin{proof}
Let $\Psi=(z\to a,z\to b)$. Its full sieves $\{z\to a\}$ and $\{z\to b\}$ and their intersection $\{z\}$ all have terminal objects, so Proposition~\ref{A:audit-prop8and9a} makes them totally $W$-aspherical. The constant functors from either arrow category to $\{z\}$ are retractions. Nevertheless the explicit ordinary final inclusion $j:\Psi\to C$ constructed and verified in Theorem~\ref{G:counterexample} has $BC\cong S^2$. It follows that $\Psi$ is neither totally $W_\infty$-aspherical nor totally $W_R$-aspherical when $R\ne0$. By Proposition~\ref{A:audit-homologygluing}, the latter localizer satisfies the stated gluing axiom. Thus the axiom does not rescue the union conclusion or either retraction version.

The failed step in the manuscript can be inspected before taking nerves. The subterminal presheaves of the two branches and their intersection are the representables $\Psi(-,a)$, $\Psi(-,b)$, and $\Psi(-,z)$. Their images under $j_!$ are the corresponding representables on $C$. Evaluating the first inclusion at $x$ gives the noninjective map
\[
 C(x,z)=\{\alpha,\beta\}\longrightarrow
 C(x,a)=\{\gamma\}.
\]
A retraction between full subcategories does not define a retraction between their subterminal presheaves: at an object outside a proper sieve, the latter would require a map from a singleton to the empty set. The source's attempted passage of categorical retractions through $j_!$ is therefore invalid, and the exhibited final image proves that its conclusion is false.
\end{proof}
